\chapter{Resultados}
\label{ch:resultados}

El objetivo de este capítulo es presentar los resultados obtenidos mediante nuestro motor numérico para la secuencia homóloga $np^{2}$, que comprende los átomos de Carbono ($Z=6$), Silicio ($Z=14$), Germanio ($Z=32$) y Estaño ($Z=50$). 
El interés de abordar esta secuencia radica en que, al transitar del carbono al estaño, la simetría angular de la subcapa de valencia se mantiene invariante. Esto nos permite aislar el problema angular y enfocar el análisis en la física radial, la cual incrementa su complejidad drásticamente por el apantallamiento impuesto por el llenado progresivo de capas internas cerradas.

\section{Validación Numérica y Teorema del Virial}

Antes de extraer cualquier parámetro espectroscópico es necesario establecer que el estado fundamental producido por el motor de B-splines es correcto. La validación descansa en dos criterios complementarios. El primero es la comparación directa de la energía total contra los valores tabulados por Froese Fischer \cite{fischer1977}, que constituyen la referencia estándar del límite Hartree-Fock no relativista. El segundo es el cociente del virial $-V/T$, que por el teorema demostrado en la Sección \ref{sec:virial} debe valer exactamente $2$ para un sistema ligado por un potencial coulombiano; su desviación cuantifica de forma global tanto la convergencia del ciclo SCF como la capacidad de la malla radial para representar simultáneamente la región nuclear y la cola asintótica de la función de onda.

El Cuadro \ref{tab:resultados_energia_global} reúne ambos criterios para los cuatro elementos de la secuencia $np^2$. El acuerdo con la referencia alcanza ocho cifras significativas en el carbono y se degrada de manera controlada al aumentar $Z$, hasta cinco cifras en el germanio. En los sistemas pesados el ciclo SCF no converge sin amortiguamiento: el apantallamiento del core reduce la separación energética entre los orbitales de valencia y los virtuales, de modo que la densidad oscila entre subespacios de una iteración a la siguiente, y fue necesario recurrir a la técnica de \textit{level shifting} descrita en la Sección \ref{sec:level_shifting}.

Conviene precisar que el core responsable de ese apantallamiento es distinto en cada caso: $[\mathrm{Ar}]\,3d^{10}\,4s^{2}$ para el germanio y $[\mathrm{Kr}]\,4d^{10}\,5s^{2}$ para el estaño. En ambos, la subcapa $d$ llena que antecede inmediatamente a la valencia es la que domina el comportamiento radial y, como veremos más adelante, la que hace indispensable la corrección de polarización del core.

\begin{table}[htbp]
    \centering
    \caption{Parámetros energéticos y cumplimiento del teorema del virial para la secuencia $np^2$, contrastados frente a los valores de referencia de Froese Fischer (F.F.) \cite{fischer1977}. La columna $\Delta$ reporta la diferencia absoluta respecto a la referencia, en Hartree para $E$ y $T$ y adimensional para el cociente.}
    \label{tab:resultados_energia_global}
    \makebox[\linewidth][c]{%
    \begin{tabular}{l l S[table-format=-4.8] S[table-format=-4.8] S[table-format=1.2e-1]}
        \toprule
        \textbf{Sistema} & \textbf{Propiedad} & {\textbf{Valor Calculado}} & {\textbf{Referencia}} & {$\Delta$} \\
        \midrule
        \multirow{3}{*}{Carbono ($Z=6$)} 
        & Energía Total $E$ & -37.68861896 & -37.688619 & 4e-8 \\
        & Cinética $T$ & 37.68861896 & 37.688619 & 4e-8 \\
        & Cociente Virial & 2.00000000 & 1.999999998 & 2e-9 \\
        \midrule
        \multirow{3}{*}{Silicio ($Z=14$)} 
        & Energía Total $E$ & -288.85435715 & -288.85436 & 2.85e-6 \\
        & Cinética $T$ & 288.85878586 & 288.85437 & 4.41e-3 \\
        & Cociente Virial & 1.99998467 & 1.999999965 & 1.53e-5 \\
        \midrule
        \multirow{3}{*}{Germanio ($Z=32$)} 
        & Energía Total $E$ & -2075.35973 & -2075.3597 & 3e-5 \\
        & Cinética $T$ & 2075.36256 & 2075.3597 & 2.86e-3 \\
        & Cociente Virial & 1.99999864 & 2.00000005 & 1.41e-6 \\
        \midrule
        \multirow{3}{*}{Estaño ($Z=50$)} 
        & Energía Total $E$ & -6022.9317 & -6022.9317 &  \\
        & Cinética $T$ & 6022.9354 & 6022.9316 & \\
        & Cociente Virial & 2.00000000 &  & 2.00000002 \\
        \bottomrule
    \end{tabular}}
\end{table}


\section{Evolución de Orbitales y Potenciales Efectivos}

Una vez superada la convergencia global de la energía, enfocamos nuestra atención en la topografía local de los orbitales de valencia. La Figura \ref{fig:wavefunctions} expone gráficamente las funciones de onda radiales $P_{np}(r)$ correspondientes a las capas de valencia activas del Carbono ($2p$), Silicio ($3p$), Germanio ($4p$) y Estaño ($5p$).

\begin{figure}[htbp]
    \centering
    \includegraphics[width=0.8\textwidth]{figures/np2_valence_wavefunctions.pdf}
    \caption{Funciones de onda radiales $P_{np}(r)$ para los electrones de valencia de la secuencia homóloga $np^2$. La normalización se mantiene y se aprecia la correcta conservación de la estructura nodal.}
    \label{fig:wavefunctions}
\end{figure}




\begin{figure}[htbp]
    \centering
    \includegraphics[width=0.8\textwidth]{figures/np2_all_orbitals.pdf}
    \caption{Distribución radial de todos los orbitales ocupados para Carbono, Silicio, Germanio y Estaño. Se observa claramente el aumento en la complejidad nodal y cómo los orbitales internos se contraen severamente hacia el origen al incrementar la carga nuclear $Z$.}
    \label{fig:all_orbitals}
\end{figure}

\begin{figure}[htbp]
    \centering
    \includegraphics[width=0.8\textwidth]{figures/np2_total_densities.pdf}
    \caption{Densidad electrónica radial total $D(r)$ del átomo para Carbono, Silicio, Germanio y Estaño. Las curvas ilustran la acumulación estriada de la densidad probabilística debida a la superposición y ocupación secuencial de las capas cerradas y la valencia.}
    \label{fig:total_densities}
\end{figure}

Como se aprecia en la Figura \ref{fig:all_orbitals}, el aumento del número cuántico principal $n$ se traduce directamente en la estructura nodal de la valencia: cero nodos para $2p$, uno para $3p$, dos para $4p$ y tres para $5p$. La gráfica superpone todas las funciones de onda radiales ocupadas del átomo neutro y hace visible el efecto del principio de exclusión, ya que conforme el \textit{core} se satura y se contrae bajo la atracción nuclear, los orbitales de valencia son desplazados hacia el exterior para mantener la ortogonalidad. La suma de estos autoestados, ponderada por su número de ocupación, produce la densidad electrónica radial total $D(r)$ de la Figura \ref{fig:total_densities}. En ella la densidad desarrolla picos pronunciados cerca del origen, dominados por las capas $1s$, $2s$ y $2p$, que se estrechan y crecen al aumentar $Z$, mientras que la cola exterior se extiende de manera progresiva al descender en el grupo. Esa cola cada vez más difusa es la manifestación de una nube electrónica más deformable frente a un campo externo, es decir, de una polarizabilidad dipolar $\alpha_d$ creciente; el efecto se vuelve cuantitativamente importante en el germanio y el estaño, y lo retomaremos al introducir el potencial de polarización de la Sección \ref{sec:core_polarization}.

La interacción multielectrónica que subyace a esta evolución se hace explícita al inspeccionar los campos medios generados. La Figura \ref{fig:potentials_eff} muestra el potencial electrostático central efectivo $V_{\text{eff}}(r)$ que percibe cada electrón de valencia $np$. Al ser un término puramente atractivo y estar dominado a distancias cortas por el campo nuclear sin apantallar $-Z/r$, el potencial diverge hacia $-\infty$ en el origen. Para comparar en una sola gráfica la divergencia del germanio con la del carbono sin perder resolución en la región de valencia, empleamos una escala pseudo-logarítmica definida por la transformación $y = \operatorname{sgn}(V) \log_{10}\left(1 + |V|/0.1\right)$, que se comporta linealmente para valores cercanos a cero y comprime de forma logarítmica las magnitudes de miles de Hartree presentes en la cercanía nuclear.

La dinámica de un electrón $p$ ($l=1$) no está regida, sin embargo, por $V_{\text{eff}}$ sino por el potencial radial efectivo $V_{\text{rad}}(r) = V_{\text{eff}}(r) + l(l+1)/2r^2$, que incorpora la repulsión centrífuga. Como se observa en la Figura \ref{fig:potentials_rad}, el término centrífugo domina cerca del origen y empuja el potencial hacia $+\infty$, de modo que es él, y no la atracción nuclear, el responsable de que la curva adopte la forma de un pozo con un mínimo bien definido.


\begin{figure}[htbp]
    \centering
    \includegraphics[width=0.8\textwidth]{figures/np2_effective_potentials.pdf}
    \caption{Potencial central electrostático puro $V_{\text{eff}}(r)$ en escala pseudo-logarítmica. Nótese cómo la interacción colapsa asintóticamente a $-\infty$ al aproximarse al origen donde el campo nuclear $-Z/r$ domina sin apantallamiento.}
    \label{fig:potentials_eff}
\end{figure}

\begin{figure}[htbp]
    \centering
    \includegraphics[width=0.8\textwidth]{figures/np2_radial_potentials.pdf}
    \caption{Potencial radial efectivo $V_{\text{rad}}(r) = V_{\text{eff}}(r) + \frac{l(l+1)}{2r^2}$ en escala lineal. La inclusión del término centrífugo ($l=1$) conforma el pozo cuántico clásico con un mínimo local bien definido.}
    \label{fig:potentials_rad}
\end{figure}

Como se aprecia en la gráfica radial, el pozo desciende desde $+\infty$ hasta un mínimo local y luego converge asintóticamente a cero. De los datos autoconsistentes, el fondo del pozo del carbono se ubica en $V_{\text{rad}} \approx -1.02 \ \text{Ha}$ a un radio $r = 0.49 \, a_0$. En el silicio, el aumento de la carga efectiva lo desplaza hasta $-15.77 \ \text{Ha}$ en $r = 0.17 \, a_0$, y en el germanio hasta $-138.88 \ \text{Ha}$ en $r = 0.07 \, a_0$.
% \todo{Falta extraer el minimo de $V_{rad}$ para el estanio y agregarlo a este parrafo.}

Conviene distinguir con claridad qué describe esta curva, porque su profundidad invita a una lectura equivocada. El potencial $V_{\text{rad}}$ gobierna exclusivamente la dinámica de \textit{un solo} electrón $np$ moviéndose en el campo promedio de los demás. El fondo de $\approx -138.88 \ \text{Ha}$ del germanio no debe confundirse con la energía potencial total del átomo que fija el teorema del virial de la Sección \ref{sec:virial}, $\ev{\hat{V}} = 2E \approx -4150 \ \text{Ha}$, la cual resulta de sumar las interacciones coulombianas de los treinta y dos electrones del sistema.

La evolución de $V_{\text{rad}}$ a lo largo de la secuencia muestra una tendencia consistente. El pozo del carbono es somero y ancho; al aumentar $Z$ el core se contrae y, pese al apantallamiento adicional que introduce, el campo coulombiano residual desplaza el mínimo hacia radios menores y profundidades mayores, hasta los $0.07 \, a_0$ del germanio. Esta contracción acerca al electrón de valencia a la región donde $\langle 1/r^3 \rangle$ crece con mayor rapidez, y constituye el mecanismo geométrico que subyace al incremento del acoplamiento espín-órbita que cuantificaremos en la sección siguiente.

Resulta instructivo confrontar estas profundidades con las distribuciones radiales de la Figura \ref{fig:all_orbitals}, porque la intuición clásica falla. Una partícula clásica reposaría en el mínimo del pozo, mientras que la cuántica se distribuye sobre toda la región permitida. Dada la marcada asimetría de $V_{\text{rad}}$, con un muro centrífugo hacia el núcleo y un decaimiento suave hacia el exterior, la función de onda se acumula preferentemente a radios mayores. Así, aunque el mínimo del carbono se sitúa en $r = 0.49 \, a_0$, el máximo de su densidad de valencia se encuentra en $r \approx 2.0 \, a_0$.

En el silicio, el germanio y el estaño este desplazamiento hacia el exterior se acentúa por el principio de exclusión. Para mantenerse ortogonal a los orbitales del core, la función de onda de valencia debe oscilar en la región de mayor atracción nuclear, lo que genera los nodos y los picos secundarios visibles cerca del origen. El resultado es que el grueso de su densidad de probabilidad queda expulsado hacia la región exterior, muy lejos del mínimo espacial del potencial que acabamos de reportar.


\section{Extracción de Parámetros de Estructura Fina}

Para calcular el espectro en acoplamiento intermedio necesitamos dos familias de parámetros, ambas extraídas directamente de la solución de campo medio. La primera son las repulsiones coulombianas dentro de la subcapa de valencia, cuantificadas por las integrales de Slater $F^k$ introducidas en la Sección \ref{sec:slater_integrals}; son las responsables de separar energéticamente los términos $^3P$, $^1D$ y $^1S$ de la configuración. La segunda es el momento radial inverso $\langle 1/r^3 \rangle$, que fija el parámetro de estructura fina $\zeta_{np}$ mediante la Ecuación \ref{eq:zeta_spin_orbit} y determina con ello el desdoblamiento de cada término en sus componentes $J$.

Los Cuadros \ref{tab:momentos_inversos} y \ref{tab:integrales_slater} contrastan los valores obtenidos por nuestro código frente a los reportados en la literatura para las cuatro especies de la secuencia.

\begin{table}[htbp]
    \centering
    \caption{Parámetros radiales críticos de estructura fina para los orbitales de valencia $np$: Momentos Inversos ($\langle 1/r^3 \rangle$) medidos en $a_0^{-3}$. V.C. denota los Valores Calculados por nuestro código, mientras que F.F. indica los valores de Froese Fischer.}
    \label{tab:momentos_inversos}
    \begin{tabular}{l c S[table-format=1.6] S[table-format=1.6] S[table-format=1.2e-1]}
        \toprule
        \textbf{Átomo} & \textbf{Capa} & {\textbf{V.C. $\langle r^{-3} \rangle$}} & {\textbf{F.F. $\langle r^{-3} \rangle$}} & {\textbf{$\Delta$}} \\
        \midrule
        Carbono ($Z=6$)   & $2p$ & 1.691811 & 1.69181 & 1e-6 \\
        Silicio ($Z=14$)  & $3p$ & 2.047231 & 2.05423 & 7.00e-3 \\
        Germanio ($Z=32$) & $4p$ & 4.793285 & 4.80120 & 7.91e-3 \\
        Estaño ($Z=50$)   & $5p$ & 6.819372 & 6.83248 & 1.31e-2 \\
        \bottomrule
    \end{tabular}
\end{table}

\begin{table}[htbp]
    \centering
    \caption{Integrales de Slater de repulsión electrostática evaluadas para la subcapa de valencia $np^2$, expresadas en Hartrees ($E_h$). V.C. indica Valores Calculados y F.F. los resultados de Froese Fischer.}
    \label{tab:integrales_slater}
    \begin{tabular}{l l S[table-format=1.8] S[table-format=1.8] S[table-format=1.2e-1]}
        \toprule
        \textbf{Átomo} & \textbf{Multipolo} & {\textbf{V.C. ($E_h$)}} & {\textbf{F.F. ($E_h$)}} & {\textbf{$\Delta$}} \\
        \midrule
        \multirow{2}{*}{Carbono} 
        & $F^0(2p, 2p)$ & 0.53860369 & 0.53860360 & 9e-8 \\
        & $F^2(2p, 2p)$ & 0.24330175 & 0.24330170 & 5e-8 \\
        \midrule
        \multirow{2}{*}{Silicio}
        & $F^0(3p, 3p)$ & 0.31650266 & 0.32971786 & 1.32e-2 \\
        & $F^2(3p, 3p)$ & 0.16272437 & 0.16593265 & 3.2e-3 \\
        \midrule
        \multirow{2}{*}{Germanio}
        & $F^0(4p, 4p)$ & 0.30137782 & 0.32624597 & 2.48e-2 \\
        & $F^2(4p, 4p)$ & 0.16272437 & 0.16593265 & 3.2e-3 \\
        \midrule
        \multirow{2}{*}{Estaño}
        & $F^0(5p, 5p)$ & 0.27874800 & 0.27875374 & 5.74e-6 \\
        & $F^2(5p, 5p)$ & 0.14723200 & 0.14722241 & 9.59e-6 \\
        \bottomrule
    \end{tabular}
\end{table}

\subsection{Escalamiento Relativista frente a la Repulsión Coulombiana}

La extracción del parámetro $\zeta_{np}$, proporcional al momento $\langle 1/r^3 \rangle$ por la Ecuación \ref{eq:zeta_spin_orbit}, junto con la repulsión coulombiana $F^2(np,np)$ definida en la Sección \ref{sec:slater_integrals}, permite localizar el punto exacto en que se rompe el acoplamiento $LS$ discutido en la Sección \ref{sec:acoplamiento_ls}. La comparación entre ambas cantidades es directa porque tienen las mismas unidades de energía y compiten por determinar la misma estructura: $F^2$ separa los términos $^3P$, $^1D$ y $^1S$ entre sí, mientras que $\zeta_{np}$ desdobla cada término en sus componentes de momento angular total $J$. El esquema $LS$ es una buena descripción mientras la primera domine a la segunda por varios órdenes de magnitud, y deja de serlo cuando ambas se vuelven comparables.

Como se derivó en la Sección \ref{sec:spin_orbita}, el parámetro de acoplamiento espín-órbita escala con la cuarta potencia de la carga nuclear efectiva, $\zeta_{nl} \sim Z_{\text{eff}}^4$, resultado de que el operador $\xi(r)$ aporta un factor $Z$ a través de la derivada del potencial y el valor esperado $\langle 1/r^3 \rangle$ aporta los tres restantes. La repulsión electrostática, en cambio, crece mucho más lentamente. Las integrales de Slater se evalúan sobre orbitales de valencia que el apantallamiento del core mantiene a un radio medio casi constante a lo largo del grupo, de manera que $F^2$ ni siquiera aumenta de forma apreciable al incrementar $Z$.


\begin{figure}[htbp]
    \centering
    \includegraphics[width=0.8\textwidth]{figures/scaling_law.pdf}
    \caption{Competencia energética entre el acoplamiento espín-órbita ($\zeta_{np}$) y la repulsión electrostática ($F^2$) a lo largo de la secuencia $np^2$, en escala doblemente logarítmica. La recta que sigue $\zeta_{np}$ corresponde a la ley de potencias $Z_{\text{eff}}^4$, mientras que $F^2$ permanece prácticamente constante; ambas escalas se aproximan a un orden de magnitud en el germanio.}
    \label{fig:scaling_law}
\end{figure}

La Figura \ref{fig:scaling_law} muestra este comportamiento en los datos de nuestro motor SCF. La repulsión electrostática $F^2$, convertida a $\text{cm}^{-1}$, decae del carbono al silicio y se estabiliza en torno a $\num{35000} \ \text{cm}^{-1}$ para el resto de la secuencia, reflejo del apantallamiento que mantiene a los orbitales de valencia en un radio medio comparable. La interacción magnética $\zeta_{np}$, en cambio, sigue una ley de potencias. Conviene precisar el sentido de esta afirmación, ya que la lectura de una recta en escala logarítmica admite dos interpretaciones distintas: en un diagrama doblemente logarítmico una dependencia $\zeta \propto Z^{p}$ aparece como una recta de pendiente $p$, mientras que un crecimiento exponencial $\zeta \propto e^{cZ}$ produciría una recta únicamente en escala semilogarítmica. La recta observada en la Figura \ref{fig:scaling_law} confirma por tanto el comportamiento polinomial con exponente $p \approx 4$ deducido en la Sección \ref{sec:spin_orbita}, y no un crecimiento exponencial.

La magnitud de la separación entre ambas escalas es lo que define el régimen de acoplamiento. En el carbono, $\zeta_{2p}$ apenas alcanza los $\sim 30 \ \text{cm}^{-1}$ frente a los $\sim \num{53000} \ \text{cm}^{-1}$ de $F^2$, de modo que el efecto magnético queda relegado a una corrección menor sobre un ordenamiento fijado por la electrostática. Al llegar al germanio, $\zeta_{4p}$ ha crecido hasta cerca de $\num{1000} \ \text{cm}^{-1}$ y la separación se ha reducido de tres órdenes de magnitud a poco más de uno. Esta diferencia entre las potencias de $Z$ con que crece cada término es la razón por la cual los átomos pesados del grupo son empujados al régimen de acoplamiento intermedio descrito en la Sección \ref{sec:acoplamiento_ls}.



La precisión alcanzada en el carbono, del orden de $10^{-8}$ Hartree respecto a la referencia del Cuadro \ref{tab:resultados_energia_global}, valida la implementación del acoplamiento angular basada en el álgebra de Racah y los operadores de Fock diferenciados de la Sección \ref{sec:rohf}. Un error en los coeficientes angulares se manifestaría como una discrepancia sistemática de magnitud mucho mayor e independiente del refinamiento de la malla, lo cual no se observa. El aumento del error relativo en el silicio, el germanio y el estaño no proviene por tanto de la parte angular, sino del crecimiento del propio sistema: cada capa cerrada adicional aporta más integrales de Slater al operador de Fock y amplía el rango dinámico radial que la malla debe cubrir, desde la región donde se concentra el orbital $1s$ hasta la cola de valencia. A ello se suman las diferencias de discretización respecto al esquema en diferencias finitas empleado por Froese Fischer.

Conviene subrayar que estos valores de referencia corresponden al límite Hartree-Fock y no al experimento. La concordancia con ellos certifica que el motor resuelve correctamente las ecuaciones que se le plantearon, pero el objetivo físico —reproducir los observables medidos— exige las correcciones de correlación y de polarización que se introducen en las secciones siguientes. La calibración del código no busca, en consecuencia, una coincidencia exacta con la referencia teórica, sino la fidelidad al límite que esa referencia define.

La conclusión central que se desprende de estos parámetros está en su escalamiento relativo. La integral $F^2(np,np)$, que gobierna la separación LS entre los términos de la configuración según la Sección \ref{sec:acoplamiento_ls}, disminuye de forma sistemática al descender en el grupo ($0.243 \to 0.163 \to 0.163 \to 0.147$ Hartree), mientras que el momento $\langle 1/r^3 \rangle$ crece por un factor cercano a cuatro en el mismo trayecto ($1.69 \to 2.05 \to 4.79 \to 6.82 \ a_0^{-3}$). Puesto que es este momento inverso el que fija $\zeta_{np}$ a través de la Ecuación \ref{eq:zeta_spin_orbit}, y con ella la dependencia $Z_{\text{eff}}^4$ obtenida en la Sección \ref{sec:spin_orbita}, ambos cuadros bastan para anticipar de manera cuantitativa la ruptura del acoplamiento puro de Russell-Saunders en la parte baja del grupo.

\subsection{Teorema de Koopmans y Energías de Ionización}

Los eigenvalores $\epsilon_{nl}$ que arroja el problema generalizado de Roothaan-Hall son, en principio, multiplicadores de Lagrange sin significado observable directo. El Teorema de Koopmans, enunciado formalmente en la Sección \ref{sec:koopmans_theorem}, les asigna uno: bajo la aproximación de orbitales congelados, el potencial de ionización del sistema queda dado por $I \approx -\epsilon_{np}$. Aunque nuestra secuencia $np^2$ es de capa abierta y el teorema se formula para capa cerrada, sigue proporcionando una estimación de primer orden útil, en el sentido perturbativo preciso que se detalla en esa sección: el error que se comete es el de la relajación de los orbitales restantes tras la remoción del electrón, contribución que aparece a segundo orden en la perturbación.

Vale la pena precisar por qué ese error no incluye una contribución de correlación de primer orden. Por el Teorema de Brillouin \cite{Szabo1989}, los elementos de matriz entre el determinante de referencia Hartree-Fock y las excitaciones simples se anulan de manera exacta, de modo que las correcciones dominantes a $-\epsilon_{np}$ provienen de la relajación orbital y de la correlación de pares, ambas de orden superior. Comparamos entonces los eigenvalores de nuestro motor SCF contra los potenciales empíricos de la base ASD del NIST, descrita en la Sección \ref{sec:nist_asd}, empleando la conversión a electrón-voltios del Apéndice \ref{ap:unidades-atomicas}. El Cuadro \ref{tab:koopmans} recoge la comparación.

\begin{table}[htbp]
    \centering
    \caption{Energías de ionización obtenidas mediante el Teorema de Koopmans ($-\epsilon_{np}$) y mediante la diferencia de energías del cálculo CI ($\Delta E$), contrastadas con los valores empíricos del NIST ASD. La última columna reporta el error relativo del resultado CI. Dado que el ión $np^1$ carece de correlación de pares en la capa de valencia, la corrección CI actúa esencialmente sobre el átomo neutro $np^2$ y por tanto aumenta el potencial de ionización predicho.}
    \label{tab:koopmans}
    \begin{tabular}{lcccc}
        \toprule
        \textbf{Átomo} & \textbf{SCF ($-\epsilon_{np}$)} & \textbf{CI ($\Delta E$)} & \textbf{NIST} & \textbf{Error (CI)} \\
        \midrule
        Carbono ($2p$) & $11.07 \text{ eV}$ & $11.09 \text{ eV}$ & $11.26 \text{ eV}$ & $1.5\%$ \\
        Silicio ($3p$) & $7.58 \text{ eV}$ & $7.71 \text{ eV}$ & $8.15 \text{ eV}$ & $5.4\%$ \\
        Germanio ($4p$) & $7.33 \text{ eV}$ & $7.40 \text{ eV}$ & $7.90 \text{ eV}$ & $6.3\%$ \\
        Estaño ($5p$) & {-} & {-} & {-} & {-} \\
        \bottomrule
    \end{tabular}
\end{table}

El Cuadro \ref{tab:koopmans} muestra que los eigenvalores de la base de splines reproducen los potenciales de ionización con un error del $1.5\%$ en el carbono, que crece hasta el $5.4\%$ en el silicio y el $6.3\%$ en el germanio. La tendencia es en sí misma informativa. El error no aumenta por una pérdida de precisión numérica, puesto que el Cuadro \ref{tab:resultados_energia_global} muestra que la energía total sigue convergiendo a cinco cifras significativas en los sistemas pesados; aumenta porque la fracción de energía que el campo medio deja de describir crece con el número de electrones. Al pasar del carbono al germanio se añaden capas internas cuya correlación con los electrones de valencia el modelo de partículas independientes ignora por construcción.

La corrección CI de valencia recupera una parte de esa energía y desplaza la predicción en la dirección correcta en los tres casos, tal como anticipa la definición de $E_{\text{corr}}$ como cantidad negativa que estabiliza el átomo neutro. Su magnitud, sin embargo, es de apenas centésimas de electrón-volt, claramente insuficiente frente a una brecha que en el germanio supera el medio electrón-volt. Esta insuficiencia es precisamente lo que motiva las dos secciones siguientes: primero el espectro multiplete, donde el mismo déficit se manifiesta de forma mucho más visible, y después la corrección de polarización del core que lo atenúa.

\subsection{Validación Experimental del Espectro Multiplete}

Los términos espectroscópicos que resultan de la configuración $np^2$ ($^3P_J$, $^1D_2$ y $^1S_0$) determinan las transiciones radiativas permitidas de cada elemento y funcionan como firmas de identificación en absorción y emisión. Conocer sus energías con precisión es lo que da utilidad práctica al cálculo: en astrofísica observacional y en física de plasmas, las líneas de los átomos neutros e ionizados constituyen el diagnóstico primario de temperatura, densidad y abundancia química del medio emisor \cite{Cowan1981}. Este es, en última instancia, el motivo que justifica el trabajo, ya que el espectro es el único punto de contacto directo entre la estructura electrónica calculada y la observación.

Para contrastar los desdoblamientos predichos en acoplamiento intermedio, extrajimos los niveles experimentales de la base ASD del NIST, descrita en la Sección \ref{sec:nist_asd}, para las especies neutras de la secuencia; en notación espectroscópica, C I, Si I, Ge I y Sn I, donde el numeral romano indica el grado de ionización. Los autoestados teóricos se obtuvieron diagonalizando el Hamiltoniano de Breit-Pauli en la base de determinantes de valencia, siguiendo el procedimiento de la Sección \ref{sec:ci_method}. Las energías resultantes, expresadas en $\text{cm}^{-1}$ y referidas al nivel $^3P_0$, se presentan en los Cuadros \ref{tab:niveles_energia_ligeros} y \ref{tab:niveles_energia_pesados}.

\begin{table}[htbp]
    \centering
    \begin{tabular}{l|rr|rr}
        \toprule
        & \multicolumn{2}{c|}{\textbf{Carbono (C I)}} & \multicolumn{2}{c}{\textbf{Silicio (Si I)}} \\
        \textbf{Término} & \textbf{HF+CI} & \textbf{NIST} & \textbf{HF+CI} & \textbf{NIST} \\
        \midrule
        $^3P_0$ & 0.0 & 0.0 & 0.0 & 0.0 \\
        $^3P_1$ & 20.8 & 16.4 & 113.6 & 77.1 \\
        $^3P_2$ & 62.2 & 43.4 & 327.1 & 223.2 \\
        $^1D_2$ & 11958.3 & 10192.6 & 8335.5 & 6298.8 \\
        $^1S_0$ & 29429.5 & 21648.0 & 17048.7 & 15394.4 \\
        \bottomrule
    \end{tabular}
    \caption{Comparación de los niveles de estructura fina ($\text{cm}^{-1}$) para los elementos más ligeros de la secuencia. Se contrasta el modelo estricto de correlación de valencia (HF+CI) contra los valores experimentales del NIST.}
    \label{tab:niveles_energia_ligeros}
\end{table}

\begin{table}[htbp]
    \centering
    \begin{tabular}{l|rrr|rrr}
        \toprule
        & \multicolumn{3}{c|}{\textbf{Germanio (Ge I)}} & \multicolumn{3}{c}{\textbf{Estaño (Sn I)}} \\
        \textbf{Término} & \textbf{HF+CI} & \textbf{CI+$V_{\text{pol}}$} & \textbf{NIST} & \textbf{HF+CI} & \textbf{CI+$V_{\text{pol}}$} & \textbf{NIST} \\
        \midrule
        $^3P_0$ & 0.0 & 0.0 & 0.0 & {-} & {-} & {-} \\
        $^3P_1$ & 869.7 & 514.5 & 557.1 & {-} & {-} & {-} \\
        $^3P_2$ & 2107.7 & 1345.5 & 1410.0 & {-} & {-} & {-} \\
        $^1D_2$ & 9648.7 & 9483.4 & 7125.3 & {-} & {-} & {-} \\
        $^1S_0$ & 19566.0 & 20265.2 & 16367.1 & {-} & {-} & {-} \\
        \bottomrule
    \end{tabular}
    \caption{Niveles de estructura fina ($\text{cm}^{-1}$) para los átomos pesados de la secuencia, referidos al nivel $^3P_0$. La columna HF+CI evidencia la sobrestimación del desdoblamiento magnético que produce el modelo de core congelado. La columna CI+$V_{\text{pol}}$ incorpora el potencial de polarización del core con $\alpha_d = 0.5$ y $r_c = 1.0\,a_0$ para el germanio.}
    \label{tab:niveles_energia_pesados}
\end{table}

\begin{figure}[htbp]
    \centering
    \includegraphics[width=1.0\textwidth]{figures/np2_energy_levels.pdf}
    \caption{Espectro de energía teórico de los términos correspondientes a la capa $np^2$. Nótese cómo la separación de los subniveles $^3P_J$ se magnifica fuertemente al descender en el grupo hacia el Germanio, rompiendo progresivamente la degeneración característica del régimen de campo débil.}
    \label{fig:np2_levels}
\end{figure}

Los Cuadros \ref{tab:niveles_energia_ligeros} y \ref{tab:niveles_energia_pesados} muestran que el modelo CI de valencia reproduce el ordenamiento jerárquico correcto de los niveles y una estimación de primer orden de los desdoblamientos, con errores que crecen de forma sistemática al descender en el grupo, como se aprecia en la Figura \ref{fig:np2_levels}.

Conviene detenerse en cómo se rompe la degeneración a lo largo de la secuencia, porque el mecanismo no cambia pero su magnitud sí. Un Hamiltoniano no relativista —el de Schrödinger para el hidrógeno o el de Hartree-Fock restringido para un átomo multielectrónico— conmuta por separado con $\mathbf{L}^2$ y $\mathbf{S}^2$ y no contiene ningún operador capaz de distinguir entre valores de $J$. En consecuencia, los niveles $^3P_0$, $^3P_1$ y $^3P_2$ emergen exactamente degenerados del ciclo SCF, y esto vale por igual para el hidrógeno y para el germanio. La degeneración no se rompe por refinar la malla ni por ampliar la base: se rompe únicamente al añadir el término $\hat{H}_{SO}$ del Hamiltoniano de Breit-Pauli, cuyo valor esperado depende de $\mathbf{l}\cdot\mathbf{s}$ y por tanto de $J$. Lo que sí cambia al bajar en el grupo es la magnitud del efecto: en el carbono el desdoblamiento es de decenas de $\text{cm}^{-1}$ y $\hat{H}_{SO}$ admite un tratamiento perturbativo sobre estados $LS$ puros, mientras que en el germanio alcanza los miles de $\text{cm}^{-1}$ y obliga a la diagonalización completa descrita en la Sección \ref{sec:acoplamiento_ls}.

El problema en el germanio es que alimentar el Hamiltoniano de Breit-Pauli con orbitales de Hartree-Fock estrictos produce un desdoblamiento demasiado grande: $\sim 870 \ \text{cm}^{-1}$ y $\sim 2100 \ \text{cm}^{-1}$ para $^3P_1$ y $^3P_2$, frente a los $557$ y $1410 \ \text{cm}^{-1}$ del NIST. La sobrestimación es de un factor cercano a $1.5$ y no puede atribuirse a la parte angular, que es la misma que reproduce correctamente el carbono. Su origen está en el momento $\langle 1/r^3 \rangle$: en un modelo de core rígido nada impide que el orbital $4p$ penetre la región interna más de lo que lo haría en el átomo real, y $\zeta_{4p}$, proporcional a ese momento por la Ecuación \ref{eq:zeta_spin_orbit}, hereda amplificada cualquier sobrestimación de la densidad cerca del núcleo.

La corrección natural consiste en admitir que el core no es rígido. El potencial de polarización $V_{\text{pol}}$, derivado en la Sección \ref{sec:core_polarization} e implementado numéricamente en la Sección \ref{sec:impl_polarizacion}, modela la respuesta dipolar del core al campo del electrón de valencia. Con los parámetros $\alpha_d = 0.5$ y $r_c = 1.0\,a_0$ para el germanio, el efecto es una relajación radial del orbital $4p$ hacia el exterior que reduce su densidad en las inmediaciones del origen y, con ella, el momento $\langle 1/r^3 \rangle$ y la constante $\zeta_{4p}$. La columna CI+$V_{\text{pol}}$ del Cuadro \ref{tab:niveles_energia_pesados} recoge el resultado: el intervalo $^3P_2$ baja de $\sim 2100$ a $\sim 1345 \ \text{cm}^{-1}$, situándose a menos del $5\%$ del valor experimental.

Que la corrección resulte necesaria en el germanio y no en los elementos ligeros tiene una explicación estructural. El core del germanio, $[\mathrm{Ar}]\,3d^{10}$, contiene una subcapa $3d$ llena que es espacialmente difusa y fácilmente polarizable; los cores del carbono ($1s^2$) y del silicio ($[\mathrm{Ne}]$) son compactos y rígidos, con polarizabilidades dipolares prácticamente nulas. Aplicar el mismo tratamiento fenomenológico en la parte alta del grupo introduciría un parámetro libre sin contenido físico, y el modelo HF+CI estricto es suficiente allí.

Es igualmente revelador que $V_{\text{pol}}$ opere casi exclusivamente sobre la separación magnética $^3P_J$ y apenas altere la separación electrostática de los singletes $^1D_2$ y $^1S_0$. La razón está en el peso radial de cada integral. La repulsión $F^2(np,np)$ se acumula sobre todo el volumen atómico y es por tanto insensible a una redistribución pequeña de la densidad, mientras que $\zeta_{np}$ depende de $\langle 1/r^3 \rangle$, cuyo integrando está fuertemente pesado hacia $r \to 0$. Una relajación leve del orbital cerca del origen basta para modificar el segundo sin alterar de forma apreciable al primero.

Conviene señalar los límites de este resultado. $V_{\text{pol}}$ es una corrección semiempírica de un solo parámetro que actúa sobre la parte radial del problema, y su éxito se concentra en los intervalos de estructura fina del término fundamental. Los términos superiores siguen mal descritos: el $^1S_0$ del germanio conserva un error superior a los $3000 \ \text{cm}^{-1}$ incluso con la corrección aplicada, y no hay motivo para esperar que un potencial central lo mejore, dado que esa desviación es de origen electrostático y no magnético. Cerrar la brecha exigiría ampliar el espacio de configuraciones o abandonar la aproximación de core congelado.

El mismo procedimiento de calibración se aplicó al estaño ($Z=50$), para el cual no disponemos de un valor de $\alpha_d$ de referencia. Se recorrió el intervalo $\alpha_d \in [2.0, 6.0]$ tomando como criterio de ajuste los niveles empíricos del NIST ($^3P_1 = 1691.8 \ \text{cm}^{-1}$ y $^3P_2 = 3427.7 \ \text{cm}^{-1}$); el perfil resultante se muestra en el Cuadro \ref{tab:estanio_ci}.

\begin{table}[htbp]
\centering
\caption{Desdoblamientos del término $^3P$ del Sn I calculados con CI en función de la polarizabilidad dipolar del core $\alpha_d$ (en unidades atómicas), con $r_c = 1.0\,a_0$. Los valores de referencia del NIST son $^3P_1 = 1691.8\ \text{cm}^{-1}$ y $^3P_2 = 3427.7\ \text{cm}^{-1}$.}
\label{tab:estanio_ci}
\begin{tabular}{cccc}
\toprule
$\alpha_d$ & $\zeta_{5p}$ (Ha) & $^3P_1$ (cm$^{-1}$) & $^3P_2$ (cm$^{-1}$) \\
\midrule
2.0 & 0.00940 & 1441.8 & 3196.7 \\
3.0 & 0.01001 & 1551.7 & 3406.6 \\
4.0 & 0.01073 & 1682.6 & 3652.7 \\
5.0 & 0.01156 & 1836.7 & 3937.5 \\
6.0 & 0.01251 & 2016.1 & 4263.4 \\
\bottomrule
\end{tabular}
\end{table}

El barrido expone una limitación de fondo del ajuste con un solo parámetro: los dos intervalos del término $^3P$ no pueden reproducirse simultáneamente. El valor $\alpha_d \approx 3.0$ acierta en el intervalo $^3P_2$ ($3406.6$ frente a $3427.7 \ \text{cm}^{-1}$) pero deja corto el $^3P_1$, mientras que $\alpha_d \approx 4.0$ ajusta el $^3P_1$ ($1682.6$ frente a $1691.8 \ \text{cm}^{-1}$) y sobreestima el $^3P_2$ en más de $200 \ \text{cm}^{-1}$. La discrepancia no es de origen radial sino angular: $V_{\text{pol}}$ es un potencial central y reescala $\zeta_{5p}$ de manera uniforme, de modo que la razón entre ambos intervalos queda fijada por los coeficientes de mezcla entre $^3P_2$ y $^1D_2$ del acoplamiento intermedio, sobre los cuales $\alpha_d$ no tiene ningún control. Adoptamos $\alpha_d \approx 3.5$ como valor de compromiso, que reparte el error entre los dos niveles.

Para complementar el resultado magnético conviene examinar el efecto cruzado de $V_{\text{pol}}$ sobre una magnitud puramente energética. El Cuadro \ref{tab:germanio_ionizacion} recoge el impacto de las correcciones sucesivas sobre el potencial de ionización del germanio.

\begin{table}[htbp]
    \centering
    \caption{Efecto de la inclusión progresiva de las correcciones sobre el potencial de ionización del germanio. La corrección CI recupera parte de la energía de correlación de valencia; $V_{\text{pol}}$ añade la contribución de la polarización del core, ausente en el modelo de core congelado. El error relativo se reduce a menos de la mitad respecto al valor de Koopmans.}
    \label{tab:germanio_ionizacion}
    \begin{tabular}{lccc}
        \toprule
        \textbf{Modelo para Ge (4p)} & \textbf{Ionización calculada} & \textbf{NIST} & \textbf{Error} \\
        \midrule
        HF puro (Koopmans) & $7.33 \text{ eV}$ & $7.90 \text{ eV}$ & $7.2\%$ \\
        HF + CI (Valencia) & $7.40 \text{ eV}$ & $7.90 \text{ eV}$ & $6.3\%$ \\
        HF + CI + $V_{\text{pol}}$ & $7.65 \text{ eV}$ & $7.90 \text{ eV}$ & $3.1\%$ \\
        \bottomrule
    \end{tabular}
\end{table}

A diferencia de $\zeta_{4p}$, que responde con fuerza a la forma exacta del orbital cerca del núcleo, la energía de ionización está determinada por la densidad en todo el dominio. Que ambas magnitudes mejoren con la misma corrección indica que $V_{\text{pol}}$ no está actuando como un mero factor de escala sobre el momento $\langle 1/r^3 \rangle$, sino que modifica el potencial efectivo de manera consistente en todo el rango radial. El error remanente del $3.1\%$ es atribuible a la correlación de valencia no capturada por el espacio CI truncado.

\section{Conclusión del Estudio $np^2$}

El estudio de la secuencia $np^2$ permite cerrar con una lectura de conjunto. El motor numérico reproduce el límite Hartree-Fock no relativista en los cuatro elementos, desde el carbono hasta el estaño, sin requerir ninguna modificación de la arquitectura de B-splines: basta ajustar la densidad de nodos y el parámetro de la malla exponencial para acomodar el rango radial creciente. Sobre esa base, la corrección CI de valencia y el potencial de polarización $V_{\text{pol}}$ pudieron incorporarse como capas sucesivas, cada una atacando un déficit identificable del modelo anterior.

El hilo que recorre el capítulo es el desplazamiento progresivo del régimen de acoplamiento. En el carbono la repulsión electrostática domina al acoplamiento magnético por tres órdenes de magnitud y el esquema $LS$ es prácticamente exacto; en el estaño esa separación se ha reducido a poco más de un orden y los términos con el mismo $J$ se mezclan de forma apreciable. El motor captura esa transición de manera cuantitativa en los intervalos de estructura fina del término fundamental, aunque los términos superiores conservan errores que las correcciones empleadas aquí no están en condiciones de eliminar.

%%% Local Variables:
%%% mode: LaTeX
%%% TeX-master: "../main"
%%% End:
